% ---------------------------------------------------------
% Section: Boolean Rings
% ---------------------------------------------------------

% ---------------------------------------------------------
% Definition: Boolean Ring
% Source: Halmos & Givant, Introduction to Boolean Algebras, \S{}1
% ---------------------------------------------------------

\begin{definitionbox}{Definition --- Boolean Ring}
\begin{definition}[Boolean Ring]
\label{def:boolean-ring}
A \textbf{Boolean ring} is an idempotent ring with unit.

That is, a Boolean ring is a ring $(R, +, \cdot, -, 0, 1)$
satisfying the additional axiom:
\[
p \cdot p = p \qquad \text{for all } p \in R. \tag{11}
\]
An element $p$ satisfying $p \cdot p = p$ is called
\textbf{idempotent}.
A ring in which every element is idempotent is called an
\textbf{idempotent ring}.

The official axiom set for a Boolean ring consists of
axioms~(1)--(3), (5), (6), (8)--(11), where axiom~(7)
$($the additive inverse law$)$ is replaced by
axiom~(10) $($characteristic~$2)$, which is derived as a theorem.
Axiom numbering follows Halmos--Givant~\cite{HalmosGivant2009}.
\end{definition}
\end{definitionbox}

\begin{remark*}[Standard quantified statement]
\[
\begin{aligned}
\operatorname{BooleanRing}(R,+,\cdot,-,0,1)
&\;\Longleftrightarrow\;
  \operatorname{Ring}(R,+,\cdot,-,0,1)
  \;\land\;
  \forall p \in R{:}\; p \cdot p = p.
\end{aligned}
\]
\end{remark*}

\begin{remark*}[Definition predicate reading]
\[
\operatorname{BooleanRing}(R,+,\cdot,-,0,1)
\;\Longleftrightarrow\;
\operatorname{Ring}(R,+,\cdot,-,0,1)
\;\land\;
\operatorname{IdempotentRing}(R,\cdot)
\]
\end{remark*}

\begin{remark*}[Negated quantified statement]
\[
\neg\,\operatorname{BooleanRing}(R,+,\cdot,-,0,1)
\;\Longleftrightarrow\;
\neg\,\operatorname{Ring}(R,+,\cdot,-,0,1)
\;\lor\;
\exists p \in R{:}\; p \cdot p \neq p.
\]
\end{remark*}

\begin{remark*}[Failure modes]
The failure modes are exactly the ways that one of the defining structural
conditions can fail.
\end{remark*}

\begin{remark*}[Failure mode decomposition]
A structure fails to be a Boolean ring in exactly one of two ways:
\[
\underbrace{\neg\,\operatorname{Ring}(R,+,\cdot,-,0,1)}_{\text{fails a ring axiom}}
\qquad\lor\qquad
\underbrace{\exists p \in R{:}\; p \cdot p \neq p}_{\text{some element is not idempotent}}
\]
The second failure mode is the structurally interesting one.
A ring in which even one element fails $p \cdot p = p$ is not Boolean.
The simplest example: the ring $\mathbb{Z}$ of integers fails
idempotence at every element except $0$ and $1$.
\end{remark*}

\begin{remark*}[Canonical examples]
\hfill
\begin{itemize}
  \item $\mathbf{2} = \{0,1\}$ with addition mod~$2$ and ordinary
        multiplication is the simplest non-degenerate Boolean ring.
  \item $\mathbf{2}^n$, the set of all $n$-tuples of elements of
        $\mathbf{2}$, with coordinatewise addition and multiplication,
        is a Boolean ring with $2^n$ elements.
        The zero is $(0,\ldots,0)$ and the unit is $(1,\ldots,1)$.
  \item $(\mathcal{P}(X), \triangle, \cap)$, the power set of any
        set $X$ with symmetric difference as addition and intersection
        as multiplication, is a Boolean ring.
        Under the correspondence $A \leftrightarrow \mathbf{1}_A$
        (the characteristic function of $A$), this ring is isomorphic
        to $\mathbf{2}^X$.
\end{itemize}
\end{remark*}

\begin{remark*}[Interpretation]
A Boolean ring is the minimal algebraic structure that forces
the full Boolean calculus to emerge from a single additional
axiom --- idempotence of multiplication.
The ring axioms supply the additive abelian group and the
multiplicative monoid with distributivity.
Idempotence then forces characteristic~$2$ (every element
equals its own additive inverse), forces commutativity of
multiplication, and identifies the ring with a field of sets
via the representation theorem (Halmos--Givant, Chapter~22).
The abstract definition thus captures exactly the algebra of
set-theoretic operations: symmetric difference plays the role
of addition and intersection plays the role of multiplication.
\end{remark*}

\begin{dependencies}
\begin{itemize}
  \item \hyperref[def:algebraic-ring]{Definition: Ring}
\end{itemize}
\end{dependencies}


% =========================================================
% Section: Boolean Rings --- Definitions and First Properties
% Chapter: Boolean Rings and Boolean Algebras --- Volume I
% Source: Halmos & Givant, Introduction to Boolean Algebras, \S{}1
% =========================================================

% ---------------------------------------------------------
% The Two-Element Boolean Ring 2
% ---------------------------------------------------------

\begin{example*}[The two-element Boolean ring]
The simplest non-degenerate Boolean ring is the two-element ring
$\mathbf{2} = \{0, 1\}$.
Its addition and multiplication are given by the following tables.

\begin{center}
\begin{tabular}{c|cc}
$+$ & $0$ & $1$ \\ \hline
$0$ & $0$ & $1$ \\
$1$ & $1$ & $0$
\end{tabular}
\qquad\qquad
\begin{tabular}{c|cc}
$\cdot$ & $0$ & $1$ \\ \hline
$0$ & $0$ & $0$ \\
$1$ & $0$ & $1$
\end{tabular}
\end{center}

\noindent
Addition in $\mathbf{2}$ is addition modulo $2$: $1 + 1 = 0$.
Multiplication in $\mathbf{2}$ is ordinary multiplication of
integers restricted to $\{0, 1\}$.
Every element of $\mathbf{2}$ satisfies $p + p = 0$ and
$p \cdot p = p$, making $\mathbf{2}$ the canonical example against
which both axioms~(10) and~(11) are verified.
\end{example*}

% ---------------------------------------------------------
% Definition: Idempotent Element
% ---------------------------------------------------------

\begin{definitionbox}{Definition --- Idempotent Element}
\begin{definition}[Idempotent Element]
\label{def:idempotent-element}
An element $p$ of a ring $R$ is called \textbf{idempotent} if
\[
p \cdot p = p.
\]
A ring in which every element is idempotent is called an
\textbf{idempotent ring}.
\end{definition}
\end{definitionbox}

\begin{remark*}[Standard quantified statement]
\[
\begin{aligned}
\operatorname{IdempotentElement}(p, R, \cdot)
&\;\Longleftrightarrow\; p \cdot p = p. \\[4pt]
\operatorname{IdempotentRing}(R, \cdot)
&\;\Longleftrightarrow\; \forall p \in R{:}\; p \cdot p = p.
\end{aligned}
\]
\end{remark*}

\begin{remark*}[Definition predicate reading]
\[
\operatorname{IdempotentElement}(p,R,\cdot)
\qquad
\operatorname{IdempotentRing}(R,\cdot)
\;\Longleftrightarrow\;
\forall p \in R{:}\;
\operatorname{IdempotentElement}(p,R,\cdot)
\]
\end{remark*}

\begin{remark*}[Negated quantified statement]
\[
\neg\,\operatorname{IdempotentElement}(p,R,\cdot)
\;\Longleftrightarrow\;
p \cdot p \neq p.
\]
\end{remark*}

\begin{remark*}[Interpretation]
Idempotence is the single additional axiom that forces Boolean
structure.  In the two-element ring $\mathbf{2}$, verification is
immediate: $0 \cdot 0 = 0$ and $1 \cdot 1 = 1$.  In a power set
algebra $\mathcal{P}(X)$, idempotence of intersection --- $A \cap A
= A$ --- is the set-theoretic counterpart.  The force of the axiom
is not in these examples but in what it implies for an arbitrary
ring: characteristic~2 and commutativity of multiplication both
follow from idempotence alone, as the theorems below establish.
\end{remark*}

\begin{dependencies}
\begin{itemize}
  \item \hyperref[def:algebraic-ring]{Definition: Ring}
\end{itemize}
\end{dependencies}

% ---------------------------------------------------------
% Definition: Characteristic 2
% ---------------------------------------------------------

\begin{definitionbox}{Definition --- Characteristic 2}
\begin{definition}[Characteristic 2]
\label{def:characteristic-2}
A ring $R$ is said to have \textbf{characteristic~$2$} if
\[
p + p = 0 \qquad \text{for all } p \in R. \tag{10}
\]
\end{definition}
\end{definitionbox}

\begin{remark*}[Standard quantified statement]
\[
\operatorname{Characteristic2}(R,+,0)
\;\Longleftrightarrow\;
\forall p \in R{:}\; p + p = 0.
\]
\end{remark*}

\begin{remark*}[Definition predicate reading]
\[
\operatorname{Characteristic2}(R,+,0)
\;\Longleftrightarrow\;
\forall p \in R{:}\;
\operatorname{IdempotentElement}(p, R, +)
\]
\end{remark*}

\begin{remark*}[Negated quantified statement]
\[
\neg\,\operatorname{Characteristic2}(R,+,0)
\;\Longleftrightarrow\;
\exists p \in R{:}\; p + p \neq 0.
\]
\end{remark*}

\begin{remark*}[Interpretation]
Characteristic~2 means that every element is its own additive
inverse: $p + p = 0$ says precisely $p = -p$.  The operation of
negation therefore becomes the identity map --- applying $-$ to
any element returns that element unchanged.  Halmos observes that
in a ring of characteristic~2, the minus sign for additive inverses
is never needed and is henceforth dropped.  Axiom~(7),
$p + (-p) = 0$, is replaced in the official Boolean ring axiom set
by the stronger equation~(10), $p + p = 0$, which is a theorem
derived from idempotence.
\end{remark*}

\begin{dependencies}
\begin{itemize}
  \item \hyperref[def:boolean-ring]{Definition: Boolean Ring}
  \item \hyperref[def:idempotent-element]{Definition: Idempotent Element}
\end{itemize}
\end{dependencies}

% ---------------------------------------------------------
% Lemma: Self Additive Inverse (Characteristic 2)
% ---------------------------------------------------------

\begin{lemmabox}{Lemma}
\begin{lemma}[Self Additive Inverse]
\label{lem:self-additive-inverse}
In any Boolean ring $R$,
\[
p + p = 0 \qquad \text{for all } p \in R.
\]

\smallskip
\noindent
\hyperref[prf:self-additive-inverse]{\textit{Go to proof.}}
\end{lemma}
\end{lemmabox}

\begin{remark*}[Standard quantified statement]
\[
\forall p \in R{:}\; p + p = 0.
\]
\end{remark*}

\begin{remark*}[Definition predicate reading]
\[
\operatorname{BooleanRing}(R,+,\cdot,-,0,1)
\;\Longrightarrow\;
\operatorname{Characteristic2}(R,+,0).
\]
\end{remark*}

\begin{remark*}[Interpretation]
This is Halmos consequence~(a): idempotence of multiplication
forces the additive structure to have characteristic~2.  The proof
expands $(p+q)^2$ by distributivity and idempotence, cancels $p$
and $q$, then sets $p = q$ to obtain $p + p = 0$.  The Lean~4
formalisation is \texttt{AdditiveIdempotence} in
\texttt{BooleanAlgebraDefinitions.lean}.
\end{remark*}

\begin{dependencies}
\begin{itemize}
  \item \hyperref[def:boolean-ring]{Definition: Boolean Ring}
  \item \hyperref[def:characteristic-2]{Definition: Characteristic 2}
\end{itemize}
\end{dependencies}

% ---------------------------------------------------------
% Lemma: Self Negation
% ---------------------------------------------------------

\begin{lemmabox}{Lemma}
\begin{lemma}[Self Negation]
\label{lem:self-negation}
In any Boolean ring $R$,
\[
-p = p \qquad \text{for all } p \in R.
\]

\smallskip
\noindent
\hyperref[prf:self-negation]{\textit{Go to proof.}}
\end{lemma}
\end{lemmabox}

\begin{remark*}[Standard quantified statement]
\[
\forall p \in R{:}\; -p = p.
\]
\end{remark*}

\begin{remark*}[Interpretation]
This follows immediately from \ref{lem:self-additive-inverse}.
Since $p + p = 0$ and $p + (-p) = 0$ both hold, the uniqueness of
additive inverses gives $-p = p$.  The negation operation is
therefore the identity map on any Boolean ring.  Halmos uses this
to drop the minus sign entirely from his notation for the rest of
the book: since $-p = p$, writing $-p$ and writing $p$ are the
same thing.
\end{remark*}

\begin{dependencies}
\begin{itemize}
  \item \hyperref[lem:self-additive-inverse]{Lemma: Self Additive Inverse}
\end{itemize}
\end{dependencies}
